
%%FILE blueprint/src/content.tex
% In this file you should put the actual content of the blueprint.
% It will be used both by the web and the print version.
% It should *not* include the \begin{document}
%
% If you want to split the blueprint content into several files then
% the current file can be a simple sequence of \input. Otherwise It
% can start with a \section or \chapter for instance.
\chapter{Introduction}

This blueprint accompanies a Lean~4 formalization of error correcting codes.
Unless stated otherwise, definitions and statements follow the conventions of
\emph{Essential Coding Theory}~\cite{GRS}, which is our main reference throughout.

\chapter{Basics of Codes}
\input{sections/Basics/definitions}
\input{sections/Basics/simple_properties}
\chapter{Linear Codes}
\input{sections/LinearCodes/definitions}
\input{sections/LinearCodes/lemmas}
\chapter{Bounds}
\input{sections/Bounds/Gilbert-Varshamov}


%%FILE blueprint/src/sections/Basics/definitions.tex
\section{Definitions}
\begin{definition}[Code]
\label{def:Code}
\lean{Code}
\leanok
A code $C$ of block length $n$ over an alphabet $\alpha$ is a subset of $\alpha^{n}$.
Typically, we will use $q$ to denote the alphabet size $|\Sigma|$.
\end{definition}

\begin{definition}[Alphabet Size]
    \label{def:q}
    \lean{Code.q}
    \leanok
    The alphabet size is $q := \lvert\alpha\rvert$, the cardinality of the finite alphabet $\alpha$.
\end{definition}

\begin{definition}[Dimension]
    \label{def:dim}
    \lean{Code.dim}
    \leanok
    \uses{def:Code, def:q}
    The dimension $k$ of a code $C$ is
    \[
        k := \log_{q}(\lvert C\rvert),
    \]
    where $q$ is the alphabet size and $\lvert C\rvert$ is the number of codewords.
\end{definition}

\begin{definition}[Rate]
    \label{def:rate}
    \lean{Code.rate}
    \leanok
    \uses{def:Code, def:dim}
    The rate $R$ of a code $C$ is
    \[
        R := \frac{k}{n}.
    \]
    where $k$ is the dimension and block length of $n$.
\end{definition}

\begin{definition}[Minimum Distance]
  \label{def:MinDim}
  \lean{Code.minDist}
  \leanok
  \uses{def:Code}
  Let $C \subseteq \Sigma^{n}$.
  The minimum distance of $C$, denoted $\Delta(C)$, is defined to be
  \[\Delta(C) = \min_{\mathbf{c_{1}} \neq \mathbf{c_{2}} \in C} \Delta(\mathbf{c_{1}}, \mathbf{c_{2}}).\]
\end{definition}

\begin{definition}[Relative Minimum Distance]
  \label{def:RelMinDim}
  \lean{Code.relMinDist}
  \leanok
  \uses{def:MinDim}
  Let $C \subseteq \Sigma^{n}$.
  The minimum distance of $C$, denoted $\delta(C)$, is defined to be
  \[\delta(C) = \min_{\mathbf{c_{1}} \neq \mathbf{c_{2}} \in C} \delta(\mathbf{c_{1}}, \mathbf{c_{2}}).\]
\end{definition}

\begin{definition}[Hamming Ball]
  \label{def:HammingBall}
  \lean{Code.hammingBall}
  \leanok
  For any vector $\mathbf{x} \in [q]^{n}$,
  \[B(\mathbf{x},e) = \{\mathbf{y} \in [q]^{n}\mid \Delta(\mathbf{x},\mathbf{y}) \leq e\}.\]
\end{definition}

\begin{definition}[Volume]
  \label{def:Vol}
  \lean{Code.hammingVolume}
  \leanok
  \uses{def:HammingBall}
  \[V_{q}(n,r) = \sum^{r}_{i = 0} \binom{n}{i}(q-1)^{i}\]
  where $r$ is the radius.
\end{definition}

\begin{definition}[Maximal With Respect to Inclusion]
  \label{def:maximalWrtInclusion}
  \lean{Code.maximalWrtInclusion}
  \leanok
  \uses{def:Code, def:q, def:MinDim}
  Let $C\subseteq \Sigma^{n}$ with minimum distance denoted $\Delta(C)$.
  $C$ is defined to be maximal with respect to inclusion if:
  \[\forall C' \subseteq \Sigma^{n}, C\subseteq C', \Delta(C) = \Delta(C') \implies C'\subseteq C\]
\end{definition}


%%FILE blueprint/src/sections/Basics/simple_properties.tex
\section{Properties}
\begin{lemma}[Dimension is at most the block length]
    \label{lem:dim_le_n}
    \lean{Code.dim_le_n}
    \leanok
    \uses{def:dim}
    For every code $C$ of block length $n$, we have $k \leq n$.
\end{lemma}

\begin{proof}
    \leanok
    \uses{def:q}
    Note that $C \subseteq \alpha^{n}$, we have $\lvert C\rvert \le \lvert\alpha\rvert^{n} = q^{n}$, thus $k = \log_{q}\lvert C\rvert \le \log_{q} q^{n} = n$.
\end{proof}
 
\begin{lemma}[Rate is at most one]
    \label{lem:rate_le_one}
    \lean{Code.rate_le_one}
    \leanok
    \uses{def:rate}
    For every code $C$, we have $R \leq 1$.
\end{lemma}
\begin{proof}
    \leanok
    \uses{lem:dim_le_n}
    By Lemma~\ref{lem:dim_le_n}, $k \leq n$, and $n \ge 0$.
    Therefore $R = k/n \leq 1$.
\end{proof}

\begin{lemma}[Sphere cardinality]
  \label{lemma:SphereCard}
  \lean{Code.ncard_hammingSphere}
  \leanok
  \uses{def:q}
  For any $\mathbf{x} \in [q]^n$ and integer $i \geq 0$,
  \[\bigl|\{\mathbf{y} \in [q]^n \mid \Delta(\mathbf{x}, \mathbf{y}) = i\}\bigr| = \binom{n}{i}(q-1)^i.\]
\end{lemma}

\begin{proof}
  \leanok
  Write $S_i = \{\mathbf{y} \in [q]^n \mid \Delta(\mathbf{x}, \mathbf{y}) = i\}$ for the sphere,
  and for a word $\mathbf{y} \in [q]^n$ let
  \[
    T(\mathbf{y}) = \{j \in [n] \mid y_j \neq x_j\}
  \]
  be its set of disagreement coordinates. By the definition of Hamming distance,
  $\Delta(\mathbf{x}, \mathbf{y}) = |\{j \in [n] \mid x_j \neq y_j\}| = |T(\mathbf{y})|$, so
  membership $\mathbf{y} \in S_i$ is the same as the condition $|T(\mathbf{y})| = i$.

  We construct an explicit bijection. Let $\mathcal{P}_i$ be the set of pairs $(T, g)$ in which $T \subseteq [n]$
  has $|T| = i$ and $g : T \to [q]$ satisfies $g(j) \neq x_j$ for every $j \in T$. Define the
  forward map
  \[
    \Phi : S_i \to \mathcal{P}_i, \qquad
    \Phi(\mathbf{y}) = \bigl(T(\mathbf{y}),\ (y_j)_{j \in T(\mathbf{y})}\bigr);
  \]
  this is well defined because $|T(\mathbf{y})| = i$ by the paragraph above, and $y_j \neq x_j$
  for each $j \in T(\mathbf{y})$ by the definition of $T(\mathbf{y})$. Define the reverse map
  \[
    \Psi : \mathcal{P}_i \to S_i, \qquad
    \Psi(T, g) = \mathbf{y}, \quad\text{where } y_j =
      \begin{cases} g(j) & \text{if } j \in T, \\ x_j & \text{if } j \notin T. \end{cases}
  \]
  Its image lies in $S_i$: the word $\mathbf{y} = \Psi(T, g)$ has $y_j \neq x_j$ exactly on $T$
  (on $T$ since $g(j) \neq x_j$, and off $T$ since $y_j = x_j$), so $T(\mathbf{y}) = T$ and
  therefore $|T(\mathbf{y})| = |T| = i$. The two maps are mutually inverse:
  \begin{itemize}
    \item $\Psi(\Phi(\mathbf{y})) = \mathbf{y}$: on $T(\mathbf{y})$ the reconstructed word takes
      the value $y_j$, and off $T(\mathbf{y})$ it takes the value $x_j$, which equals $y_j$
      because $j \notin T(\mathbf{y})$ means exactly that $y_j = x_j$.
    \item $\Phi(\Psi(T, g)) = (T, g)$: we showed $T(\Psi(T, g)) = T$, and on $T$ the values kept
      by $\Phi$ are the $g(j)$.
  \end{itemize}
  Therefore, $|S_i| = |\mathcal{P}_i|$.

  It remains to count $\mathcal{P}_i$. Grouping the pairs by their first coordinate $T$,
  \[
    |\mathcal{P}_i| = \sum_{\substack{T \subseteq [n] \\ |T| = i}}
      \bigl|\{\, g : T \to [q] \mid g(j) \neq x_j \text{ for all } j \in T \,\}\bigr|.
  \]
  For a fixed $T$, specifying such a $g$ means choosing, independently for each $j \in T$, one of
  the $q - 1$ values in $[q] \setminus \{x_j\}$; there are thus $(q-1)^{|T|} = (q-1)^i$ such
  functions. As the number of $i$-element subsets of $[n]$ is $\binom{n}{i}$,
  \[
    |S_i| = |\mathcal{P}_i|
      = \sum_{\substack{T \subseteq [n] \\ |T| = i}} (q-1)^i
      = \binom{n}{i}(q-1)^i.
  \]
\end{proof}

\begin{lemma}[Ball cardinality equals volume]
  \label{lemma:BallCard}
  \lean{Code.ncard_hammingBall}
  \leanok
  \uses{def:HammingBall, def:Vol, def:q}
  For any $\mathbf{x} \in [q]^n$ and $r \geq 0$,
  \[|B(\mathbf{x}, r)| = \Vol_q(n, r).\]
\end{lemma}

\begin{proof}
  \leanok
  \uses{lemma:SphereCard}
  For each integer $0 \leq i \leq r$ let $S_i = \{\mathbf{y} \in [q]^n \mid \Delta(\mathbf{x}, \mathbf{y}) = i\}$
  be the sphere of radius $i$ about $\mathbf{x}$.

  The $S_i$ are pairwise disjoint: if $\mathbf{y} \in S_i \cap S_j$ then
  $i = \Delta(\mathbf{x}, \mathbf{y}) = j$, so distinct indices give disjoint spheres, as a word
  has a single Hamming distance to $\mathbf{x}$.

  The $S_i$ also cover the ball. By the definition of the Hamming ball,
  $\mathbf{y} \in B(\mathbf{x}, r)$ means $\Delta(\mathbf{x}, \mathbf{y}) \leq r$. Since
  $\Delta(\mathbf{x}, \mathbf{y})$ is a nonnegative integer, this holds iff
  $\Delta(\mathbf{x}, \mathbf{y}) = i$ for exactly one $i$ with $0 \leq i \leq r$, i.e. iff
  $\mathbf{y} \in S_i$ for some such $i$. Therefore,
  \[
    B(\mathbf{x}, r) = \bigsqcup_{i=0}^{r} S_i,
  \]
  a disjoint union by the previous paragraph.

  Finally, we compute the cardinalities. Because the union is disjoint, its cardinality is the sum of the
  cardinalities of the parts, and each part is counted by Lemma~\ref{lemma:SphereCard}:
  \[
    |B(\mathbf{x}, r)|
      = \sum_{i=0}^{r} |S_i|
      = \sum_{i=0}^{r} \binom{n}{i}(q-1)^i.
  \]
  The right-hand side is exactly $\Vol_q(n, r)$ by the definition of the volume, so
  $|B(\mathbf{x}, r)| = \Vol_q(n, r)$.
\end{proof}

\begin{lemma}[Volume entropy bound]
\label{lemma:VolEntropy}
\lean{Code.hammingVolume_le_pow_mul_entropy}
\leanok
\uses{def:q, def:Vol}
Let $q \geq 2$ and $n \geq 1$ be integers, and let $r$ be an integer with $0 \leq r$ and $\frac{r}{n} \leq 1 - \frac{1}{q}$. Then
\[
  \Vol_q(r, n) \leq q^{\,n H_q(r/n)}.
\]
\end{lemma}
\begin{proof}
  Write $\lambda = \frac{r}{n}$. If $r = 0$ then $\Vol_q(0, n) = 1 = q^{n H_q(0)}$, so assume $r \geq 1$, giving $0 < \lambda \leq 1 - \frac{1}{q} < 1$. From the definition of $H_q$,
  \[
    q^{n H_q(\lambda)} = (q-1)^{r}\, \lambda^{-r}\,(1-\lambda)^{-(n-r)},
  \]
  so it suffices to show $\Vol_q(r,n)\,(q-1)^{-r}\lambda^{r}(1-\lambda)^{n-r} \leq 1$. Put $\theta = \frac{\lambda}{(q-1)(1-\lambda)}$. Note that $\lambda \leq 1 - \frac{1}{q}$ gives $\frac{\lambda}{1-\lambda} \leq q-1$, hence $0 \leq \theta \leq 1$. For each $0 \leq i \leq r$, since $r - i \geq 0$ and $\theta \leq 1$,
  \[
    \binom{n}{i} (q-1)^{i} (q-1)^{-r}\lambda^{r}(1-\lambda)^{n-r} = \binom{n}{i} \lambda^{i} (1-\lambda)^{n-i}\,\theta^{\,r-i} \leq \binom{n}{i} \lambda^{i} (1-\lambda)^{n-i}.
  \]
  Summing over $0 \leq i \leq r$ and extending the sum to $0 \leq i \leq n$,
  \[
    \Vol_q(r,n)\,(q-1)^{-r}\lambda^{r}(1-\lambda)^{n-r} \leq \sum_{i=0}^{n}\binom{n}{i} \lambda^{i}(1-\lambda)^{n-i} = (\lambda + (1-\lambda))^{n} = 1,
  \]
  the last step by the binomial theorem. Rearranging gives $\Vol_q(r,n) \leq q^{n H_q(\lambda)}$.
\end{proof}


%%FILE blueprint/src/sections/Bounds/Gilbert-Varshamov.tex
\begin{lemma}[Covering]
\label{lemma:covers}
\lean{Code.covers}
\leanok
\uses{def:Code, def:q, def:MinDim, def:Vol, def:maximalWrtInclusion, def:HammingBall}
Let $q \geq 2$ and $n \geq d \geq 1$ be integers, and let $C \subseteq [q]^{n}$ be a
code of minimum distance at least $d$ that is maximal with respect to inclusion.
For $\mathbf{c} \in C$, let
\[
  B(\mathbf{c}, d-1) = \{ \mathbf{z} \in [q]^{n} : \Delta(\mathbf{c}, \mathbf{z}) \le d-1 \}
\]
denote the Hamming ball of radius $d-1$ centered at $\mathbf{c}$.
Then
\[
  \bigcup_{\mathbf{c} \in C} B(\mathbf{c}, d-1) = [q]^{n}.
\]
\end{lemma}

\begin{proof}
  Let $\mathbf{z} \in [q]^{n}$ and suppose towards a contradiction that $\mathbf{z} \notin B(c, d-1)$ for every $\mathbf{c} \in C$, that is $\Delta(\mathbf{c}, \mathbf{z}) \geq d$ for all $\mathbf{c} \in C$.
  Note that $\mathbf{z} \neq \mathbf{c}$ for every $\mathbf{c}$ since $d \geq 1$ and thus $\mathbf{z} \notin C$.
  Then $C \cup \{\mathbf{z}\}$ is a code where any two distinct codewords lie at distance greater than $d$ meaning that $C$ is not maximal which is a contradiction.
  Hence every $\mathbf{z}$ lies in $B(c,d-1)$ and
  \[[q]^{n} = \bigcup_{\mathbf{c} \in C} B(\mathbf{c},d-1).\]
\end{proof}

\begin{lemma}[Maximal Packing]
\label{lemma:MaximalPacking}
\lean{Code.maxPacking}
\leanok
\uses{def:Code, def:q, def:MinDim, def:Vol, def:maximalWrtInclusion}
Let $q \geq 2$ and $n \geq d \geq 1$ be integers, and let $C \subseteq [q]^n$ be
maximal, with respect to inclusion, among all codes of minimum distance at
least $d$.
Then
\[
  |C| \geq \frac{q^n}{\Vol_q(n,d-1)}.
\]
\end{lemma}

\begin{proof}
  \leanok
  \uses{lemma:covers, lemma:BallCard}
  For every codeword $c \in C$ we define a ball $B(\mathbf{c},d-1)$.
  Note that by Lemma~\ref{lemma:covers}, the union of these balls cover $[q]^{n}$.
  Then we must have the following inequality,
  \[
  q^{n} =
  |[q]^{n}| = 
  \left|\bigcup_{\mathbf{c} \in C}B(\mathbf{c}, d-1)\right| \leq
  \sum_{\mathbf{c} \in C} |B(\mathbf{c},d-1)| =
  |C| \cdot \Vol_{q}(n,d-1) \implies
  |C| \geq \frac{q^{n}}{\Vol_{q}(n,d-1)}.
  \]
\end{proof}

\begin{theorem}[Gilbert-Varshamov Bound; general codes]
\label{thm:GVBgc}
\lean{Code.gilbert_varshamov}
\leanok
\uses{def:Code, def:q, def:rate, def:RelMinDim, lemma:MaximalPacking, lemma:VolEntropy}
See~\cite{GRS}.
Let $q \geq 2$ be an integer.
For every $0 \leq \delta < 1 - \frac{1}{q}$ and every positive integer $n$, there exists a $q$-ary code of block length $n$ with rate $R \geq 1 - H_q(\delta)$ and relative distance at least $\delta$.
\end{theorem}

\begin{proof}
    \leanok
    Fix $\delta$ with $0 \leq \delta < 1 - \frac{1}{q}$ and a positive integer $n$, and put
    \[
      d = \max(\lceil \delta n \rceil, 1).
    \]
    Then $d \geq 1$, and $d \leq n$ since $\delta < 1$ gives $\lceil \delta n \rceil \leq n$, so $n \geq d \geq 1$. Note also $d \geq \delta n$ and $d - 1 \leq \delta n$.
    The codes $C \subseteq [q]^{n}$ of minimum distance at least $d$ form a finite nonempty family (any singleton works), so we may pick one, call it $C$, that is maximal with respect to inclusion.
    By Lemma~\ref{lemma:MaximalPacking},
    \[
      |C| \geq \frac{q^{n}}{\Vol_{q}(n,d-1)}.
    \]
    Since $\Vol_{q}(n,d-1) \leq q^{n}$ we have $|C| \geq 1$, so taking $\log_{q}$ and dividing by $n$,
    \[
      R = \frac{\log_{q}|C|}{n} \geq 1 - \frac{\log_{q}\Vol_{q}(d-1,n)}{n}.
    \]
    Note that $0 \leq \frac{d-1}{n} \leq \delta < 1 - \frac{1}{q}$, so by Lemma~\ref{lemma:VolEntropy} at $\lambda = \frac{d-1}{n}$,
    \[
      \Vol_{q}(n,d-1) \leq q^{\,n H_q(\frac{d-1}{n})} \implies \frac{\log_{q}\Vol_{q}(d-1,n)}{n} \leq H_q\!\left(\tfrac{d-1}{n}\right),
    \]
    and hence $R \geq 1 - H_q(\frac{d-1}{n})$. Finally $H_q$ is nondecreasing on $[0, 1 - \frac{1}{q}]$ by the monotonicity of $H_q$ and $\frac{d-1}{n} \leq \delta$, so $H_q(\frac{d-1}{n}) \leq H_q(\delta)$ and
    \[
      R \geq 1 - H_q(\delta).
    \]
    The relative distance of $C$ is at least $\frac{d}{n} \geq \delta$.
    Hence $C$ is a $q$-ary code of block length $n$ with rate $R \geq 1 - H_q(\delta)$ and relative distance at least $\delta$ as desired.
\end{proof}

\begin{theorem}[Gilbert-Varshamov Bound; linear codes]
\uses{def:Code, def:q, def:rate, def:MinDim, def:LinearCode, def:GeneratorMatrix}
See~\cite{GRS}.
Let $q$ be a prime power.
For every $0 \leq \delta < 1 - \frac{1}{q}$, every $0 \leq \varepsilon \leq 1 - H_q(\delta)$, and every positive integer $n$, set $k = n\bigl(1 - H_q(\delta) - \varepsilon\bigr)$.
If $G \in \mathbb{F}^{k \times n}_{q}$ is picked uniformly at random, then $G$ generates a code of rate $1 - H_q(\delta) - \varepsilon$ and relative distance at least $\delta$ with probability strictly greater than $1 - q^{-\varepsilon n}$.
\end{theorem}

\begin{proof}
    TODO
\end{proof}


%%FILE blueprint/src/sections/LinearCodes/definitions.tex
\section{Definitions}
\begin{definition}[Linear Code by Subspace]
\label{def:LC_subspace}
\lean{LinearCode}
\leanok
\uses{def:Code}
Let $q$ be a prime power (i.e. $q = p^{s}$ for some prime $p$ and integer $s \geq 1$).
$C \subseteq \mathbb{F}^{n}_{q}$ is a linear code if it is a linear subspace of $\mathbb{F}^{n}_{q}$.
If $C$ has dimenension $k$ and distance $d$ then it will be referred to as an $[n,k,d]_{q}$ or just an $[n,k]_{q}$ code.\footnote{Our lean declaration does not directly depend on code.}
\end{definition}

\begin{definition}[Generator Matrix]
\label{def:GeneratorMatrix}
\lean{LinearCode.GeneratorMatrix}
\leanok
\uses{def:LC_subspace}
If $C$ is an $[n, k]_{q}$ linear code then there exists a matrix $G \in \mathbb{F}^{k \times n}_{q}$ of rank $k$ satisfying
\[C = \{x \cdot G \mid x \in \mathbb{F}^{k}_{q}.\}\]
We define $G$ as a generator matrix of $C$.
Note that $G$ is not unique.
\end{definition}

\begin{definition}[Parity Check Matrix]
\label{def:ParityCheckMatrix}
\lean{LinearCode.ParityCheckMatrix}
\leanok
\uses{def:LC_subspace}
If $C$ is an $[n, k]_{q}$ linear code then there exists a matrix $H \in \mathbb{F}^{(n-k) \times n}_{q}$ of rank $n - k$ satisfying
\[C = \{\mathbf{y} \in \mathbb{F}^{n}_{q} \mid H \cdot \mathbf{y}^{\intercal} = 0\}.\]
We define $H$ as a partiy check matrix of $C$.
Note that $H$ is not unique.
\end{definition}


%%FILE blueprint/src/sections/LinearCodes/lemmas.tex
\section{Lemmas}

\begin{lemma}[Subspace Implies Generator]
    \uses{def:LC_subspace,def:GeneratorMatrix}
    \label{lemma:subspace-imp-generator}
    Let $C$ be a linear code that is a $k$-dimensional subspace of $\mathbb{F}^{n}_{q}$.
    Then there exists a $k \times n$ matrix $G$ with rank $k$ over $\mathbb{F}_q$ such that
    \[
        C = \{\mathbf{f}G \mid \mathbf{f} \in \mathbb{F}^{k}_{q}\}.
    \]
\end{lemma}

\begin{proof}
Note that $C$ is a $k$-dimensional subspace of $\mathbb{F}_q^n$, and thus has a basis $\{\mathbf{g}_1, \ldots, \mathbf{g}_k\}$.
Let $G$ be the $k \times n$ matrix whose $i$-th row is $\mathbf{g}_i$.
For any $\mathbf{f} = (f_1, \ldots, f_k) \in \mathbb{F}_q^k$, the row-vector product is
\[
    \mathbf{f}G = \sum_{i=1}^k f_i \mathbf{g}_i.
\]
Thus $\{\mathbf{f}G \mid f \in \mathbb{F}^{k}_{q}\}$ is the set of $\mathbb{F}_q$-linear combinations of $\mathbf{g}_{1},\ldots, \mathbf{g}_{k}$. 
Since $\mathbf{g}_{1}, \ldots, \mathbf{g}_{k}$ is a basis of $C$ and the span equals $C$, we must have $C = \{\mathbf{f}G \mid f \in \mathbb{F}^{k}_{q}\}$. It also follows that $rank(G)=k$ since $G$'s rows are basis vectors.
\end{proof}

\begin{lemma}[Generator Implies Subspace]
    \uses{def:LC_subspace, def:GeneratorMatrix}
    \label{lemma:generator-imp-subspace}
    Let $G \in \mathbb{F}^{k \times n}_{q}$ with rank $k$ and let $C = \{\mathbf{f}G \mid \mathbf{f} \in \mathbb{F}^{k}_{q}\}$.
    Then $C$ is a $k$-dimensional linear code over $\mathbb{F}^{n}_{q}$.
\end{lemma}

\begin{proof}
    First of all, 
    \begin{itemize}
        \item $\forall \mathbf{c}\in C$, $G$ constrains $\mathbf{c}$ to be $n$ dimensional
        \item A combination of: 
        \begin{itemize}
            \item The closure of $\mathbb{F}_q$ under addition and multiplication 
            \item The fact of $\mathbf{c}=\mathbf{f}G$ for some $f\in \mathbb{F}^{k}_{q}$ 
        \end{itemize}
        means all elements of $\mathbf{c}$ are in $\mathbb{F}_q$
    \end{itemize}
    This means $\mathbf{c}\in \mathbb{F}^{n}_{q}$ and therefore $C\subseteq \mathbb{F}^{n}_{q}$. 
    Given this, it is sufficient to prove $C$ contains the zero vector and is closed under linear addition and scalar multiplication.
    \begin{itemize}
        \item We can show that the zero vector $\mathbf{0}\in \mathbb{F}^{n}_{q}$ is in $C$ because $\mathbf{0}G=\mathbf{0}$.
        \item We can show that $C$ is closed under addition by considering that $\forall \mathbf{a},\mathbf{b}\in C$,
            $\exists \mathbf{x},\mathbf{y}\in \mathbb{F}^{k}_{q}$ such that $\mathbf{a}=\mathbf{x}G,\mathbf{b}=\mathbf{y}G$,
            distributivity of matrix multiplication gives us $\mathbf{a}+\mathbf{b}=\mathbf{x}G+\mathbf{y}G=(\mathbf{x}+\mathbf{y})G$,
            and since $\mathbf{x}+\mathbf{y}\in \mathbb{F}^{k}_{q}$ because of closure under addition,
            $\mathbf{a}+\mathbf{b}=(\mathbf{x}+\mathbf{y})G\in C $
        \item We can show that $C$ is closed under scalar multiplication by considering that $\forall \mathbf{a}\in C$ that
            $\exists \mathbf{x}\in \mathbb{F}^{k}_{q}$ such that $\mathbf{a}=\mathbf{x}G$.
            And therefore, $\forall d\in \mathbb{F}_{q}, d\mathbf{a}=d\mathbf{x}G=(d\mathbf{x})G$.
            And since $d\mathbf{x}\in \mathbb{F}^{k}_{q}$ by its own closure under scalar multiplication, it must follow $d(\mathbf{x}G)\in C$
        These three conditions should be sufficient to prove that $C$ is a valid vector space.
    \end{itemize}
    Therefore $C$ is a subspace of $\mathbb{F}^{n}_{q}$, and because $\{\mathbf{f}G \mid \mathbf{f} \in \mathbb{F}^{k}_{q}\}$ is the set of $\mathbb{F}_q$-linear combinations of linearly independent $\mathbf{g}_{1},\ldots, \mathbf{g}_{k}$, it follows that $dim(C)=k$.
\end{proof}

\begin{lemma}[Subspace implies Parity Check]
    \uses{def:LC_subspace,def:ParityCheckMatrix,lemma:subspace-imp-generator}
    \label{lemma:subspace-imp-parity-check}
    Let $C$ be a linear code that is a $k$-dimensional subspace of $\mathbb{F}^{n}_{q}$.
    Then there exists a matrix $H \in \mathbb{F}^{(n-k) \times n}_{q}$ with rank $n-k$ such that
    \[
        C = \{\mathbf{f} \in \mathbb{F}^{n}_{q} \mid H\mathbf{f}^{\intercal} = \mathbf{0}\}.
    \]
\end{lemma}

\begin{proof}
    Suppose we take the orthogonal complement of $C$, $C^{\perp}\subseteq \mathbb{F}^{n}_{q}$. This gives us that
    \begin{itemize}
        \item $dim(C^{\perp})=n-dim(C)=n-k$
        \item From Lemma~\ref{lemma:subspace-imp-generator}, we know that we can get some $H\in \mathbb{F}^{n-k \times n}_{q}$ such that
    \[C^{\perp} = \{\mathbf{f}H \mid \mathbf{f} \in \mathbb{F}^{n-k}_{q}\}\]
        \item And that from Lemma~\ref{lemma:subspace-imp-generator} the rows of $H$ are in $C^{\perp}$
        \item From Lemma~\ref{lemma:subspace-imp-generator}, we get generator $G$ for $C$, where $G$'s rows are all basis elements and therefore elements of $C$
    \end{itemize}
    This means that $HG^T$ is the zero matrix since every element is the dot product of some element of $C$ and $C^{\perp}$. This gives us that
    \[ \forall \mathbf{f}\in \mathbb{F}^k_q, HG^T\mathbf{f}^T = H(\mathbf{f}G)^T = \mathbf{0}\]
    And by our definition of the generator matrix,
    \[\forall \mathbf{c} \in C, \exists \mathbf{f}\in \mathbb{F}^k_q, \mathbf{f}G = \mathbf{c}\]
    \[\implies \forall \mathbf{c} \in C, H\mathbf{c}^T = H(\mathbf{f}G)^T = \mathbf{0}\]
    \[\implies C \subseteq \{\mathbf{c} \in \mathbb{F}^{n}_{q} \mid H\mathbf{c}^{\intercal} = \mathbf{0}\}\]
    To make notation easier, let's say this subspace (we know it's a subspace because of Lemma~\ref{lemma:parity-check-imp-subspace} below) is $C'=\{\mathbf{c} \in \mathbb{F}^{n}_{q} \mid H\mathbf{c}^{\intercal} = \mathbf{0}\}$. We know that
    \begin{itemize}
        \item $dim(\mathbb{F}^{n}_{q})=n$
        \item $dim(C)=k$
    \end{itemize}
    We can assert that the rank of $H$ is $n-k$ because Lemma~\ref{lemma:subspace-imp-generator} stipulates that $H$ has full rank. From here we can use the rank-nullity theorem to say that, since $C'$ represents the null space of $H$
    \[dim(C')+rank(H)=dim(\mathbb{F}^{n}_{q})\]
    \[\implies dim(C') + n-k = n\]
    \[\implies dim(C') = k\]
    and since $dim(C')=dim(C)$ and $C\subseteq C'$, it should follow that $C=C'$.
\end{proof}

\begin{lemma}[Parity Check Implies subspace]
    \uses{def:LC_subspace, def:ParityCheckMatrix}
    \label{lemma:parity-check-imp-subspace}
    Let $H \in \mathbb{F}^{(n-k) \times n}_{q}$ with rank $n-k$ and let $C = \{\mathbf{f} \in \mathbb{F}^{n}_{q} \mid H\mathbf{f}^{\intercal} = \mathbf{0}\}$.
    Then $C$ is a subspace of $\mathbb{F}^{n}_{q}$ with dimension $k$.
\end{lemma}

\begin{proof}
Because every $\mathbf{c}\in C$ is defined as being in $\mathbb{F}^{n}_{q}$, $C \subseteq \mathbb{F}^{n}_{q}$. It should then be sufficient to prove that $C$ contains the zero vector and is closed under linear addition and scalar multiplication
\begin{itemize}
    \item The zero vector $\mathbf{0}\in \mathbb{F}^{n-k}_{q}$ is in $C$ because $H\mathbf{0}^T=\mathbf{0}$
    \item We can show that $C$ is closed under addition by considering that $\forall \mathbf{a},\mathbf{b}\in C, H\mathbf{a}^T=H\mathbf{b}^T=\mathbf{0}$, meaning $H\mathbf{a}^T+H\mathbf{b}^T= \mathbf{0}+\mathbf{0} = \mathbf{0}$
        and therefore by distributivity of matrix multiplication and distributivity of transverses, We have that
        \[H\mathbf{a}^T+H\mathbf{b}^T = H(\mathbf{a}^T+\mathbf{b}^T) = H(\mathbf{a}+\mathbf{b})^T = \mathbf{0}\]
        and since $(\mathbf{a}+\mathbf{b})\in \mathbb{F}^{n}_{q}$ by closure of $\mathbb{F}^{n}_{q}$ under addition, $H(\mathbf{a}+\mathbf{b})^T = \mathbf{0}$, and $\mathbf{a}+\mathbf{b}\in C$
    \item We can show that $C$ is closed under scalar multiplication by considering that $\forall \mathbf{a}\in C$ that
        $H\mathbf{a}^T=0$, then $\forall d\in \mathbb{F}_{q}$, that $H\mathbf{a}^T=\mathbf{0} \implies dH\mathbf{a}^T = d*0=\mathbf{0}$
        which $\implies H(d\mathbf{a})^T=\mathbf{0}$ because we can commute scalars with matrices and vectors,
        giving us that $d\mathbf{a}\in C$
\end{itemize}
These three conditions prove that $C$ is a vector space, and is the kernel of $H$. Using rank-nullity theorem, we can see:
\[rank(H)+dim(C)=dim(\mathbb{F}^{n}_{q})\]
\[\implies n-k+dim(C)=n\]
\[\implies dim(C)=k\]

\end{proof}

